\documentclass[a4paper, serif, 10pt]{moderncv}

%% moderncv
% color and style
\moderncvcolor{black}
\moderncvstyle{banking}

%% page layout/margins
\usepackage[top=2.5cm, bottom=2.5cm, left=1.5cm, right=1.5cm]{geometry}

\nopagenumbers{}

%% Babel config for Spanish language
% \usepackage[spanish]{babel} has some conflicting issues with moderncv
% so we have to use enable the following options to make it work
%
% ref:
% - https://tex.stackexchange.com/a/140161/36007
\usepackage[spanish,es-lcroman]{babel}

%% fontspec
\usepackage{fontspec}

\IfFontExistsTF{Linux Libertine}{
  \setmainfont[Ligatures={TeX, Common}, Numbers=OldStyle]{Linux Libertine}
}{}
\IfFontExistsTF{Linux Libertine O}{
  \setmainfont[Ligatures={TeX, Common}, Numbers=OldStyle]{Linux Libertine O}
}{}

%% CTeX
% CJK support, used to show CJK characters in the resume
%
% - heading=false: disable ctex heading
% - punct=kaiming: use kaiming punctuations styles for CJK
% - scheme=plain: use plain scheme, do not override \normalsize font size
% - space=auto: space settings for CJK characters
%
% CJK fonts are resolved per locale language, e.g. Noto Serif/Sans CJK JP
% for Japanese resumes, with fallback to the next region if not installed
%
% ref:
% - http://ctan.mirrorcatalogs.com/language/chinese/ctex/ctex.pdf
\usepackage[UTF8, heading=false, punct=kaiming, scheme=plain, space=auto]{ctex}

\IfFontExistsTF{Noto Serif CJK SC}{
  \setCJKmainfont{Noto Serif CJK SC}
}{}
\IfFontExistsTF{Noto Sans CJK SC}{
  \setCJKsansfont{Noto Sans CJK SC}
}{}

%% line spacing
\usepackage{setspace}
\setstretch{1.125}

%% URL styling - use normal text instead of monospace
\urlstyle{same}

%% Patch \httplink and \httpslink to handle full URLs
% The original moderncv macros always prepend http:// or https:// to URLs.
% This causes issues when the URL already has a protocol (e.g., https://example.com)
% resulting in malformed URLs like https://https://example.com.
% This patch checks if the URL already contains "://" and uses it directly if so.
% Note: We use \str_set:Nx (x-expansion) to expand macros like \@homepage first.
\ExplSyntaxOn
\str_new:N \g_yamlresume_url_str

\RenewDocumentCommand{\httpslink}{O{}m}{
  \str_set:Nx \g_yamlresume_url_str {#2}
  \str_if_in:NnTF \g_yamlresume_url_str {://}
    { \url{#2} }
    { \url{https://#2} }
}

\RenewDocumentCommand{\httplink}{O{}m}{
  \str_set:Nx \g_yamlresume_url_str {#2}
  \str_if_in:NnTF \g_yamlresume_url_str {://}
    { \url{#2} }
    { \url{http://#2} }
}
\ExplSyntaxOff

\name{Alejandro Vega Moreno}{}
\title{Ingeniero de IA especializado en visión por computador, NLP y MLOps}
\phone[mobile]{+34 612 345 678}
\email{alejandro.vega.moreno@gmail.com}
\homepage{https://alejandrovega.dev}

\address{Calle de Serrano 45, 4º B, Madrid, Comunidad de Madrid, España, 28001}{}{}

\extrainfo{{\small \faLinkedin}\ \href{https://www.linkedin.com/in/alejandrovega-ai}{@alejandrovega-ai} {} {} {} • {} {} {} 
{\small \faGithub}\ \href{https://github.com/alejandrovega-ai}{@alejandrovega-ai} {} {} {} • {} {} {} 
{\small \faStackOverflow}\ \href{https://stackoverflow.com/users/1234567/alejandro-vega}{@alejandro-vega}}

\begin{document}

\maketitle

\section{Información básica}

\cvline{}{Ingeniero de IA con más de 6 años de experiencia diseñando, desplegando y monitorizando sistemas de aprendizaje automático en producción. Especializado en visión por computador, procesamiento de lenguaje natural y MLOps, con impacto medible en eficiencia operativa y experiencias de cliente.
%
\begin{itemize}
\item Lideré el ciclo completo de modelos de IA, desde la exploración de datos hasta el despliegue en cloud y la monitorización continua.
\item Experiencia con Python, PyTorch, TensorFlow, Kubernetes y plataformas de datos en AWS y Azure.
\item Colaborador técnico y mentor, con comunicación clara entre equipos de negocio, producto e ingeniería.
\end{itemize}}
\section{Educación}

\cventry{sept 2017–jul 2019}
        {Maestría, Inteligencia Artificial, Puntuación: 9,2/10}
        {Universidad Politécnica de Madrid}
        {\url{https://www.upm.es}}
        {}
        {\begin{itemize}
\item Formación orientada a la investigación aplicada y al desarrollo de sistemas inteligentes en entornos reales.
\item Trabajo final centrado en detección de anomalías con modelos híbridos y despliegue en cloud.
\end{itemize}
\textbf{Cursos}: Aprendizaje Automático Avanzado, Visión por Computador, Procesamiento de Lenguaje Natural, Sistemas de Recomendación, Computación Distribuida, Ética y Gobernanza de la IA}

\cventry{sept 2013–jun 2017}
        {Licenciatura, Ingeniería Informática, Puntuación: 8,7/10}
        {Universidad de Sevilla}
        {\url{https://www.us.es}}
        {}
        {\begin{itemize}
\item Base sólida en programación, bases de datos, estadística y desarrollo de software.
\item Participación en proyectos académicos de análisis de datos y automatización de procesos.
\end{itemize}
\textbf{Cursos}: Algoritmos y Estructuras de Datos, Bases de Datos, Ingeniería del Software, Estadística, Sistemas Operativos, Redes de Computadores}
\section{Experiencia laboral}

\cventry{mar 2022 hasta la fecha}
        {Ingeniero de IA Senior}
        {Nexo IA Labs}
        {\url{https://www.nexoia.es}}
        {}
        {\begin{itemize}
\item Diseñé y desplegué modelos de visión por computador y NLP para productos de análisis documental y atención al cliente.
\item Lideré la implantación de pipelines MLOps con MLflow, Docker y Kubernetes, reduciendo el tiempo de puesta en producción de modelos en un 40\%.
\item Definí métricas de calidad, deriva de datos y monitorización de modelos, mejorando la fiabilidad de los sistemas en producción.
\item Colaboré con producto y negocio para priorizar casos de uso de IA según impacto y viabilidad técnica.
\item Mentoricé a ingenieros junior y promoví buenas prácticas de experimentación, revisión de código y documentación.
\end{itemize}
\textbf{Palabras clave}: Visión por computador, MLOps, PyTorch, AWS, Liderazgo técnico}

\cventry{sept 2019–feb 2022}
        {Ingeniero de Aprendizaje Automático}
        {DataMente Soluciones}
        {\url{https://www.datamente.es}}
        {}
        {\begin{itemize}
\item Desarrollé modelos predictivos para segmentación de clientes, previsión de demanda y detección de fraude.
\item Construí APIs con FastAPI y servicios de inferencia escalables sobre AWS, integrados con plataformas internas de datos.
\item Automaticé procesos de entrenamiento y validación con Airflow, reduciendo tareas manuales repetitivas.
\item Realicé análisis exploratorios y presentaciones ejecutivas para equipos no técnicos, traduciendo resultados en decisiones de negocio.
\end{itemize}
\textbf{Palabras clave}: Machine Learning, FastAPI, AWS, Airflow, Analítica predictiva}
\section{Idiomas}

\cvline{Español}{Competencia nativa o bilingüe \hfill \textbf{Palabras clave}: Nativo}
\cvline{Inglés}{Competencia profesional plena \hfill \textbf{Palabras clave}: TOEIC 940, C1 acreditado}
\cvline{Catalán}{Competencia profesional plena \hfill \textbf{Palabras clave}: Nivel C2, Uso profesional}
\section{Competencias}

\cvline{Aprendizaje automático}{Experto \hfill \textbf{Palabras clave}: Python, PyTorch, TensorFlow, scikit-learn, XGBoost}
\cvline{Procesamiento de lenguaje}{Avanzado \hfill \textbf{Palabras clave}: Transformers, BERT, LLMs, spaCy, RAG}
\cvline{Visión por computador}{Avanzado \hfill \textbf{Palabras clave}: OpenCV, YOLO, Segmentación, CNN, TensorRT}
\cvline{MLOps e infraestructura}{Avanzado \hfill \textbf{Palabras clave}: Docker, Kubernetes, MLflow, Airflow, CI/CD}
\cvline{Datos y cloud}{Avanzado \hfill \textbf{Palabras clave}: AWS SageMaker, Azure ML, Spark, SQL, Snowflake}
\cvline{Desarrollo de software}{Avanzado \hfill \textbf{Palabras clave}: Python, FastAPI, Git, C++, Pruebas unitarias}
\section{Premios}

\cventry{may 2024}
        {Nexo IA Labs}
        {Premio a la Innovación en IA}
        {}
        {}
        {Reconocimiento interno por el diseño de un sistema de visión artificial que redujo los errores de clasificación documental en un 35\%.}

\cventry{jul 2019}
        {Universidad Politécnica de Madrid}
        {Mejor proyecto de fin de máster}
        {}
        {}
        {Galardón al proyecto de detección temprana de anomalías en series temporales mediante modelos híbridos y despliegue en cloud.}
\section{Certificados}

\cventry{abr 2023}
        {Amazon Web Services}
        {AWS Certified Machine Learning - Specialty}
        {\url{https://aws.amazon.com/certification/certified-machine-learning-specialty/}}
        {}
        {}

\cventry{nov 2021}
        {Scrum.org}
        {Professional Scrum Master I}
        {\url{https://www.scrum.org/professional-scrum-certifications/professional-scrum-master-i-certification}}
        {}
        {}
\section{Publicaciones}

\cventry{oct 2022}
        {Detección temprana de anomalías en series temporales con modelos híbridos}
        {Revista Iberoamericana de Inteligencia Artificial}
        {\url{https://revista.iberoai.org/articulos/deteccion-anomalias-series-temporales}}
        {}
        {\begin{itemize}
\item Propone una combinación de modelos estadísticos y redes neuronales para detectar anomalías en datos industriales.
\item Evalúa el enfoque en conjuntos de datos reales y analiza su robustez ante ruido y cambios de distribución.
\end{itemize}}
\section{Proyectos}

\cventry{feb 2023–nov 2023}
        {Asistente conversacional para consultas sobre documentación clínica mediante recuperación aumentada.}
        {Asistente clínico con RAG}
        {\url{https://github.com/alejandrovega-ai/asistente-clinico-rag}}
        {}
        {\begin{itemize}
\item Implementé un pipeline de ingestión, indexación vectorial y recuperación de documentos médicos.
\item Añadí evaluación de respuestas y trazas para reducir alucinaciones y mejorar la trazabilidad.
\item Desplegué el servicio con FastAPI, Docker y una base vectorial gestionada en cloud.
\end{itemize}
\textbf{Palabras clave}: RAG, LLMs, FastAPI, Búsqueda vectorial}

\cventry{may 2021–ene 2022}
        {Sistema de inspección visual automática para detectar defectos en línea de producción.}
        {Visión artificial para control de calidad}
        {\url{https://github.com/alejandrovega-ai/vision-calidad}}
        {}
        {\begin{itemize}
\item Entrené modelos de detección y segmentación con imágenes industriales etiquetadas por expertos.
\item Optimicé la inferencia con TensorRT para cumplir requisitos de latencia en planta.
\item Diseñé un panel de revisión humana para validar casos ambiguos y retroalimentar el modelo.
\end{itemize}
\textbf{Palabras clave}: Visión por computador, YOLO, TensorRT, Control de calidad}
\section{Intereses}

\cvline{Deporte}{Escalada, Ciclismo, Senderismo}
\cvline{Música}{Guitarra, Jazz, Producción musical}
\cvline{Divulgación tecnológica}{Meetups, Mentoría, Artículos técnicos}
\section{Voluntariado}

\cventry{ene 2021–dic 2024}
        {Mentor y ponente}
        {Comunidad IA España}
        {\url{https://www.comunidadia.es}}
        {}
        {\begin{itemize}
\item Mentoricé a profesionales en transición hacia roles de IA y aprendizaje automático.
\item Impartí charlas sobre MLOps, evaluación de modelos y despliegue responsable de IA.
\item Colaboré en la organización de encuentros técnicos y talleres prácticos para la comunidad.
\end{itemize}}
    
\end{document}